\subsection{逆极限非空的充分条件}

在本小节中，我们将给出逆极限非空的两个最常用的充分条件（另见习题 5）。

命题5. 设 \((\mathbf{E}_{\alpha}, f_{\alpha \beta})\) 是相对于一个有可数共尾子集的有向集 I 的集合逆系统，并且进一步假设诸 \(f_{\alpha \beta}\) 是满射。那么，如果 \(\mathbf{E} = \lim \mathbf{E}_{\alpha}\) ，则典范映射 \(f_{x}: \mathbf{E} \to \mathbf{E}_{\alpha}\) 对每个 \(\alpha \in \mathbf{I}\) 都是满射（并且，特别地， \(\mathbf{E}\) 在诸 \(\mathbf{E}_{\alpha}\) 均非空的前提下非空）。

让 \((\alpha_{n})\) 为I的元素的一个序列，它构成I的一个共尾子集。由于I是有向的，我们可以归纳地由条件定义I的元素的一个序列 \((\beta_{n})\) ： \(\beta_{0} = \alpha_{0}, \beta_{n} \geqslant \beta_{i}\) 对于 \(i < n\) 并且 \(\beta_{n} \geqslant \alpha_{n}\) 。显然该序列 \((\beta_{n})\) 是递增的，并构成 I 的共尾子集。根据第 1 号的命题 1 及关系式 \(f_{\alpha} = f_{\alpha \beta_{n}} \circ f_{\beta_{n}}\) （对于 \(\alpha \leqslant \beta_{n}\) ），我们只需对 I = N 的情形证明该命题。此外，显然只需证明 \(f_{0}\) 是满射。设 \(x_{0} \in \mathrm{E}_{0}\) 。定义 \(x_{n} \in \mathrm{E}_{n}\) （对于 \(n \geqslant 1\) ）归纳地取为集合 \(\overline{f}_{n-1,n}(x_{n-1})\) 的一个元素，这是可行的，因为后者集合根据假设是非空的。然后我们对 \(n - m\) 进行归纳，证明 \(x_{m} = f_{mn}(x_{n})\) （对于 \(m \leqslant n\) ），由此得出 \(x = (x_{n})\) 属于E。

第二个准则涉及相对于指标集I的逆系统 \((\mathbf{E}_{\alpha}, f_{\alpha\beta})\) ，使得对每个 \(\alpha \in I\) 我们给定 \(E_{\alpha}\) 的子集的一个集合 \(S_{\alpha}\) ，满足以下条件：

\begin{enumerate}
\def\labelenumi{(\roman{enumi})}
\tightlist
\item
  属于 \(\mathfrak{S}_{\alpha}\) 的集合的每个交集也属于 \(\mathfrak{S}_{\alpha}\) 。
\end{enumerate}

特别地（通过考虑空族的交集）可以得出 \(\mathbf{E}_{\alpha} \in \mathfrak{S}_{\alpha}\) .

\begin{enumerate}
\def\labelenumi{(\roman{enumi})}
\setcounter{enumi}{1}
\tightlist
\item
  如果子集族 \(\mathfrak{F} \subset \mathfrak{S}_{\alpha}\) 使得属于 \(\mathfrak{F}\) 的集合的每个有限交集是非空的，那么 \(\bigcap_{\mathbf{M} \in \mathfrak{F}} \mathbf{M}\) 是非空的
\end{enumerate}

鉴于(i)，显然(ii)等价于以下条件：

(ii)' 若 \(\mathfrak{G} \subset \mathfrak{S}_{\alpha}\) 是左有向的（关于包含关系）且不包含空集，则 \(\bigcap_{M \in \mathfrak{G}} M\) 是非空的

定理1。假设I是有向的，且这些集合 \(\mathfrak{S}_{\alpha}\) 满足条件(i)和(ii)，且逆系统 \((\mathrm{E}_{\alpha}, f_{\alpha\beta})\) 具有以下性质：

\begin{enumerate}
\def\labelenumi{(\roman{enumi})}
\setcounter{enumi}{2}
\item
  对于每一对指标 \(\alpha, \beta\) 使得 \(\alpha \leqslant \beta\) ，以及每一个 \(x_{\alpha} \in \mathbf{E}_{\alpha}\) ，我们有 \(\overline{f}_{\alpha\beta}(x_{\alpha}) \in \mathfrak{S}_{\beta}\) .
\item
  对于每一对指标 \(\alpha, \beta\) 使得 \(\alpha \leqslant \beta\) ，以及每一个 \(\mathbf{M}_{\beta} \in \mathfrak{S}_{\beta}\) ，我们有 \(f_{\alpha \beta}(\mathbf{M}_{\beta}) \in \mathfrak{S}_{\alpha}\) .
\end{enumerate}

让 \(\mathbf{E} = \varprojlim \mathbf{E}_{\alpha}\) 并令 \(f_{\alpha}: \mathbf{E} \to \mathbf{E}_{\alpha}\) 为每个 \(\alpha \in \mathbf{I}\) 的典范映射。则：

\begin{enumerate}
\def\labelenumi{(\alph{enumi})}
\tightlist
\item
  对于每一个 \(\alpha \in I\) 我们拥有
\end{enumerate}

\[
f _ {\alpha} (\mathbf {E}) = \bigcap_ {\beta \geqslant \alpha} f _ {\alpha \beta} (\mathbf {E} _ {\beta}).\tag{19}
\]

\begin{enumerate}
\def\labelenumi{(\alph{enumi})}
\setcounter{enumi}{1}
\tightlist
\item
  如果 \(\mathbf{E}_{\alpha}\) 对每个 \(\alpha \in \mathbf{I}\) 都是非空的，那么 \(\mathbf{E}\) 是非空的
\end{enumerate}

让 \(\Sigma\) 为由满足以下条件的族 \(\mathfrak{A} = (\mathrm{A}_{\alpha})_{\alpha \in \mathbf{I}}\) 所组成的集合：

\begin{enumerate}
\def\labelenumi{(\arabic{enumi})}
\setcounter{enumi}{19}
\tightlist
\item
\end{enumerate}

\[
\mathrm{A} _ {\alpha} \neq \emptyset \quad \text{且} \quad \mathrm{A} _ {\alpha} \in \mathfrak {S} _ {\alpha} \quad \text{对所有} \quad \alpha \in \mathrm{I};\tag{21}
\]

\[
f _ {\alpha \beta} (\mathbf {A} _ {\beta}) \subset \mathbf {A} _ {\alpha} \quad \text{当} \quad \alpha \leqslant \beta .
\]

如果 \(\mathfrak{A} = (\mathrm{A}_{\alpha})\) 和 \(\mathfrak{A}' = (\mathrm{A}_{\alpha}')\) 是属于 \(\Sigma\) 的任意两个元素，设关系 \(\mathfrak{A} \leqslant \mathfrak{A}'\) 表示 \(\mathrm{A}_{\alpha} \supset \mathrm{A}_{\alpha}'\) 对于所有 \(\alpha\) 。显然 \(\Sigma\) 按此关系是有序的。

\begin{enumerate}
\def\labelenumi{(\arabic{enumi})}
\item
  让我们首先证明有序集 \(\Sigma\) 是归纳的。设 \(L\) 是全序集，且设 \(\lambda \to \mathfrak{A}^{\lambda} = (A_{\alpha}^{\lambda})_{\alpha \in I}\) 是从 \(L\) 到 \(\Sigma\) 中的严格递增映射。对于每一个 \(\alpha \in I\) ，令 \(B_{\alpha} = \bigcap_{\lambda \in L} A_{\alpha}^{\lambda}\) 。那么立即可以看出族 \(\mathfrak{B} = (B_{\alpha})_{\alpha \in I}\) 满足(21)；由(i)和(ii)可知，它也满足(20)，因此属于 \(\Sigma\) ；并且显然 \(\mathfrak{B}\) 是诸 \(\mathfrak{A}^{\lambda}\) 所成集合的上界。
\item
  让 \(\mathfrak{A} = (\mathrm{A}_{\alpha})\) 是 \(\Sigma\) 的一个极大元。我们将证明 \(\mathrm{A}_{\alpha} = f_{\alpha \beta}(\mathrm{A}_{\beta})\) 每当 \(\alpha \leqslant \beta\) 。 设 \(\mathrm{A}_{\alpha}' = \bigcap_{\beta \geqslant \alpha} f_{\alpha \beta}(\mathrm{A}_{\beta})\) 对于所有 \(\alpha \in \mathbf{I}\) ，且让我们证明 \(\mathfrak{A}' = (\mathrm{A}_{\alpha}')\) 属于 \(\Sigma\) 。首先注意，若 \(\alpha \leqslant \beta \leqslant \gamma\) ，我们有 \(f_{\alpha \gamma}(\mathrm{A}_{\gamma}) = f_{\alpha \beta}(f_{\beta \gamma}(\mathrm{A}_{\gamma})) \subset f_{\alpha \beta}(\mathrm{A}_{\beta})\) 由（21）；此外， \(f_{\alpha \beta}(\mathrm{A}_{\beta}) \in \mathfrak{S}_{\alpha}\) 由（iv），且 \(f_{\alpha \beta}(\mathrm{A}_{\beta}) \neq \emptyset\) 由（20）。因此条件（i）和（ii）表明 \(\mathfrak{A}'\) 满足（20）。另外，若 \(\alpha \leqslant \beta\) ，我们有
\end{enumerate}

\[
f _ {\alpha \beta} (\mathbf {A} _ {3} ^ {\prime}) \subset \bigcap_ {\gamma \geqslant \beta} f _ {\alpha \beta} (f _ {\beta \gamma} (\mathbf {A} _ {\gamma})) = \bigcap_ {\gamma \geqslant \beta} f _ {\alpha \gamma} (\mathbf {A} _ {\gamma});
\]

且对每个 \(\delta \geqslant \alpha\) 存在 \(\gamma \in I\) 使得 \(\gamma \geqslant \delta\) 并且 \(\gamma \geqslant \beta\) ，因此 \(f_{\alpha \gamma}(A_{\gamma}) \subset f_{\alpha \delta}(A_{\delta})\) ，从而

\[
\bigcap_ {\gamma \geqslant \beta} f _ {\alpha \gamma} (\mathrm{A} _ {\gamma}) = \bigcap_ {\delta \geqslant \alpha} f _ {\alpha \delta} (\mathrm{A} _ {\delta}) = \mathrm{A} _ {\alpha} ^ {\prime}.
\]

因此， \(\mathfrak{A}'\) 满足（21），因此属于 \(\Sigma\) 。由于 \(\mathbf{A}_{\alpha}' \subset \mathbf{A}_{\alpha}\) 对于所有 \(\alpha\) ，由 \(\mathfrak{A}\) 在 \(\Sigma\) 中的极大性蕴含 \(\mathfrak{A}' = \mathfrak{A}\) ，于是我们的断言得证。

\begin{enumerate}
\def\labelenumi{(\arabic{enumi})}
\setcounter{enumi}{2}
\tightlist
\item
  接下来我们将证明，若 \(\mathfrak{A} = (\mathbf{A}_{\alpha})\) 是 \(\Sigma\) 的一个极大元素，则每个 \(\mathbf{A}_{\alpha}\) 都由单个元素组成。设 \(x_{\alpha} \in \mathbf{A}_{\alpha}\) 。对于每一个 \(\beta \geqslant \alpha\) ，令 \(\mathbf{B}_{\beta} = \mathbf{A}_{\beta} \cap f_{\alpha \beta}(x_{\alpha})\) ；若 \(\beta\) 不 \(\geqslant \alpha\) ，令 \(\mathbf{B}_{\beta} = \mathbf{A}_{\beta}\) ，则我们将看到 \(\mathfrak{B} = (\mathbf{B}_{\beta})\) 属于 \(\Sigma\) 。如果 \(\beta\) 不是 \(\geqslant \alpha\) ，关系 \(\beta \leqslant \gamma\) 蕴含 \(f_{\beta \gamma}(\mathbf{B}_{\gamma}) \subset f_{\beta \gamma}(\mathbf{A}_{\gamma}) \subset \mathbf{A}_{\beta} = \mathbf{B}_{\beta}\) 。另一方面，若 \(\alpha \leqslant \beta \leqslant \gamma\) ，则由于
\end{enumerate}

\[
\overline {{f}} _ {\alpha \gamma} ^ {- 1} (x _ {\alpha}) = \overline {{f}} _ {\beta \gamma} ^ {- 1} (\overline {{f}} _ {\alpha \beta} ^ {- 1} (x _ {\alpha})),
\]

我们拥有

\[
f _ {\beta \gamma} (\overline {{f _ {\alpha \gamma} ^ {- 1}}} (x _ {\alpha})) \subset \overline {{f _ {\alpha \beta} ^ {- 1}}} (x _ {\alpha}),
\]

并且由于 \(f_{\beta \gamma}(\mathbf{A}_{\gamma}) \subset \mathbf{A}_{\beta}\) ，我们再次有 \(f_{\beta \gamma}(\mathbf{B}_{\gamma}) \subset \mathbf{B}_{\beta}\) ，因此族 \(\mathfrak{B}\) 满足 (21)。由于 \(\mathbf{A}_{\alpha} = f_{\alpha \beta}(\mathbf{A}_{\beta})\) 每当 \(\alpha \leqslant \beta\) 由证明的第(2)部分可知，显然 \(B_{\beta} \neq \varnothing\) 对于所有 \(\beta \in I\) 。最后，根据 (i) 和 (iii)，我们有 \(B_{\beta} \in S_{\beta}\) 对于所有 \(\beta \in I\) ，因此 \(B \in \Sigma\) 。由于 \(B_{\beta} \subset A_{\beta}\) 对于所有 \(\beta \in I\) ，由 A 的极大性蕴含 \(B_{\beta} = A_{\beta}\) 对于所有 \(\beta\) ，特别地 \(A_{\alpha} = \{x_{\alpha}\}\) .

\begin{enumerate}
\def\labelenumi{(\arabic{enumi})}
\setcounter{enumi}{3}
\tightlist
\item
  我们现在可以证明定理 1。让我们从 (a) 开始。我们有
\end{enumerate}

\[
f _ {\alpha} (\mathbf {E}) \subset \bigcap_ {\beta \geqslant \alpha} f _ {\alpha \beta} (\mathbf {E} _ {\beta}).
\]

反之，设

\[
x _ {\alpha} \in \bigcap_ {\beta \geqslant x} f _ {\alpha \beta} (\mathrm{E} _ {\beta}),
\]

并置

\[
\mathbf {B} _ {\beta} = \overline {{f}} _ {\alpha \beta} ^ {1} (x _ {\alpha})
\]

若 \(\beta \geqslant \alpha\) 则如上式，否则置 \(B_{\beta} = E_{\beta}\) 。根据 \(x_{\alpha}\) 的定义，诸 \(B_{\beta}\) 是非空的，并且我们有 \(B_{\beta} \in \mathfrak{S}_{\beta}\) 对于所有 \(\beta \in I\) 根据(iii)和(i)，此外，显然 \(f_{\beta \gamma}(B_{\gamma}) \subset B_{\beta}\) 每当 \(\beta \leqslant \gamma\) 。因此 \(\mathfrak{B} = (B_{\beta}) \in \Sigma\) . 设 \(\mathfrak{A} = (A_{\alpha})\) 是 \(\Sigma\) 的一个极大元，使得 \(\mathfrak{A} \geqslant \mathfrak{B}\) （\(\mathfrak{A}\) 的存在性由(1)及§2第4号定理2的推论1得出)。由于根据(3)， \(A_{\beta}\) 具有形式 \(\{y_{\beta}\}\) 对于所有 \(\beta \in I\) ，由此可得 \(y = (y_{\beta})\) 属于 \(E\) ，以及 \(f_{\alpha}(y) = y_{\alpha} = x_{\alpha}\) 根据定义。

最后，(a)蕴含(b)。我们可以不妨假设I非空（否则无需证明）。假设诸 \(E_{\alpha}\) 为非空意味着 \(f_{\alpha\beta}(E_{\beta}) \neq \emptyset\) 对于所有 \(\beta \geqslant \alpha\) 。由于诸 \(f_{\alpha\beta}(E_{\beta})\) （对于固定的 \(\alpha\) 和可变的 \(\beta \geqslant \alpha\) ）构成由属于 \(S_{\alpha}\) 的、 \(E_{\alpha}\) 的子集组成的一个左有向集，条件(ii)'证明了

\[
\bigcap_ {\beta \geqslant \alpha} f _ {\alpha \beta} (\mathbf {E} _ {\beta}) \neq \varnothing .
\]

因此， \(f_{\alpha}(\mathbf{E}) \neq \emptyset\) 由(a)可知，进而 \(\mathbf{E} \neq \emptyset\) .

证毕。

注记。假设定理1中的条件(iii)被替换为以下较弱条件：

(iii)' 对于每个 \(\alpha \in I\) 和每个非空集合 \(M_{\alpha} \in \mathfrak{S}_{\alpha}\) 存在 \(x_{\alpha} \in M_{\alpha}\) 使得 \(f_{\alpha \beta}(x_{\alpha}) \in \mathfrak{S}_{\beta}\) 对于每一个 \(\beta \geqslant \alpha\) .

则定理1的结论(b)仍然成立；定理1的第(1)部分和第(2)部分的证明保持不变，第(3)部分的证明仍然

有效，前提是我们小心地取 \(x_{\alpha} \in \mathbf{A}_{\alpha}\) 使得 \(f_{\alpha \beta}(x_{\alpha})\) 属于 \(\mathfrak{S}_{\beta}\) 对于所有 \(\beta \geqslant \alpha\) 。最后，(4)的证明表明，如果

\[
\bigcap_ {\beta \geqslant \alpha} f _ {\alpha \beta} (\mathbf {E} _ {\beta}) \neq \varnothing
\]

并且如果我们在该集合中选择一个 \(x_{\alpha}\) ，使得 \(\overline{f_{\alpha\beta}^{-1}(x_{\alpha})} \in \mathfrak{S}_{\beta}\) 每当 \(\beta \geqslant \alpha\) , 就存在 \(y \in E\) 使得 \(f_{\alpha}(y) = x_{\alpha}\) ，这证明了我们的断言。

\minorheading{示例}

\begin{enumerate}
\def\labelenumi{(\arabic{enumi})}
\tightlist
\item
  如果 \(E_{\alpha}\) 是有限集，则可以通过取 \(S_{\alpha}\) 为 \(E_{\alpha}\) 的所有子集的集合来应用定理1。* 这个例子在一般拓扑学中被推广到 \(E_{\alpha}\) 是紧拓扑空间的情形， \(f_{\alpha\beta}\) 连续映射，以及 \(S_{\alpha}\) 为 \(E_{\alpha}\) 的闭子集的集合（一般拓扑学，第一章，第9节，第6小节)。*
\end{enumerate}

* (2) 设 A 是一个具有单位元的环，并且对每个 \(\alpha \in I\) 让 \(T_{\alpha}\) 为一个Artinian左A-模。令 \(E_{\alpha}\) 是 \(T_{\alpha}\) 的一个齐性空间， \(T_{\alpha}\) 忠实地作用于其上（使得 \(E_{\alpha}\) 是附属于 \(T_{\alpha}\) 的仿射空间）。对于 \(\beta \geqslant \alpha\) ，假设 \(f_{\alpha \beta}: E_{\beta} \to E_{\alpha}\) 是一个仿射映射。取 \(\mathfrak{S}_{\alpha}\) 为由空集和 \(E_{\alpha}\) 中的仿射线性簇组成的集合。则条件(i)平凡地满足，而(ii)由以下事实得出： \(T_{\alpha}\) 是阿廷的；因为这蕴含在诸集合 \(M \in \mathfrak{F}\) 的有限交集所成的集合中存在一个极小元素，且该极小元素必须等于 \(\bigcap_{M \in \mathfrak{F}} M\) 。最后，由于 \(f_{\alpha \beta}\) 是仿射的，条件（iii）和（iv）显然成立。*

